\section{Results and Analysis}
\label{sec:results}

This chapter makes one argument in five moves, and the order is load-bearing. The field's standard
metric is shown to be gameable before any model is scored, because every later claim depends on the
ruler. Under a ruler that cannot be gamed, the model is drug-blind, and the blindness is located in
what the model is asked to predict rather than in what it can represent. Re-encoding the prediction
target then produces genuine prompt sensitivity --- and the chapter's central result is that this
sensitivity does not become reliable held-out prediction: a training-only lookup table beats the
model decisively, and the gap between the two is where the thesis ends.

\subsection{The metric is the first result}
\label{sec:res-metrics}

Before asking what the model knows, we have to know what the scoreboard rewards. The answer, for
the field's default metric, is: nothing about the drug. A predictor that places every gene at the
middle rank --- carrying no information about drug, cell or biology --- scores $\DEdr(K{=}50) =
0.9999$ [$0.99991$, $0.99992$], above the two-real-cell noise ceiling ($0.76$) and above the model
($0.73$). The ordering is inverted: less information yields a higher score, which is the signature
of an artefact rather than of skill. The same predictor's partial-$\DEdr$ and $\paneltau$ are
\emph{undefined} rather than merely poor --- its shift from control is identically zero, so there
is nothing left to correlate once the control is regressed out. The cause is exactly the mechanism
Section~\ref{sec:methods-metrics} flagged: $\DEdr$ selects the scored genes using the truth, and a
prediction that reverts toward the centre correlates with the true change by construction.

\begin{table}[h]
\centering
\begin{tabular}{lrrr}
\toprule
Predictor & Information & $\DEdr$ (K50) & partial-$\DEdr$ \\
\midrule
\texttt{revert\_center} (all genes $\rightarrow P/2$) & none & \textbf{0.9999} & undefined \\
\texttt{revert\_mean} (predict mean control) & none, no fit & 0.961 & 0.25 \\
ridge linear (control $\rightarrow$ shift) & control & 0.947 & 0.28 \\
noise ceiling (two real cells) & --- & 0.76 & --- \\
model (1B LLM) & control + drug & 0.73 & not measured \\
\bottomrule
\end{tabular}
\caption{$\DEdr$ is saturated by a control artefact. Regressing the control out of the two fitted
baselines leaves $\approx 0.26$ of real but \emph{drug-agnostic} skill, and a baseline that performs
no fitting at all reaches the same figure. The model's own partial-$\DEdr$ was never computed --- it
requires a GPU pass that was not run --- so the cell is marked rather than filled with the expected
value; the model's $\paneltau$, which \emph{is} measured, carries the same conclusion below.}
\label{tab:de-artifact}
\end{table}

Removing the control (partial-$\DEdr$) collapses the raw $0.95$ to $\approx 0.26$ for the two
baselines that were scored --- and that residue is matched by a baseline that performs no fitting at
all. The model's own partial-$\DEdr$ is not measured, so the comparison is anchored instead to
$\paneltau$, which is measured for every arm and tells the same story: \texttt{revert\_mean} $0.268$,
ridge $0.265$, model $0.257$ (tier 2, \texttt{worst} convention). The response decomposes into
control reversion (an artefact), a generic drug program ($\approx 0.26$, drug-agnostic, trivially
matched), and a drug-specific component that no predictor captures.

Grading the five metrics against each other settles which instrument the rest of the thesis may
use. Under the hardened \DRF{} protocol of Section~\ref{sec:methods-metrics} --- tie-aware ranks,
the ratio resampled whole, the family controlled by Holm's procedure --- $\NIR$ is the only
calibrated metric: $\DRF = \ci{+0.635}{+0.618}{+0.652}$, Holm $p = 0.0025$. The four rivals remain
strongly negative and none approaches significance (Holm-adjusted $p = 1.0000$ for all four); in
the earlier unhardened calibration they read $-0.081$ (\texttt{weighted\_r2}), $-0.231$
($\paneltau$), $-0.415$ (\texttt{spearman\_expr}) and $-0.448$ ($\DEdr$), and a negative \DRF{}
means the metric scores a drug-agnostic baseline \emph{above} the noise ceiling. The inversion is a
property of the data, not a leak: a leave-one-out correction barely moves the numbers, since real
cell lines carry 20--204 drugs, and it survives holding the plate constant.

This replicates Miller et al.\ --- and then diverges from them. Their central methodological claim
is confirmed emphatically: metric choice dominates, the rank-based metric is the one that survives,
and the correlation-based metrics they identify as poorly calibrated are exactly the ones that fail
here. The agreement extends to the specific metrics --- they find $\NIR$ consistently calibrated
across 14 Perturb-seq datasets, and we find it is the only calibrated metric of the five we tested
on drug perturbation. Our \texttt{weighted\_r2} result is also consistent with their finding that
weighted metrics behave better than unweighted ones. The divergence comes at the next step. Miller
et al.\ conclude that under calibrated metrics deep models outperform uninformative baselines;
Section~\ref{sec:res-blind} evaluates on the calibrated metric and finds the model at chance. Two
differences in setting are candidates for the discrepancy --- genetic versus drug perturbation, and
a discriminative regression model versus an autoregressive generative one --- and
Section~\ref{sec:res-lookup} supplies a third and more structural possibility.

\begin{figure}[htbp]
\centering
\includegraphics[width=\textwidth]{fig-calibration}
\caption[The standard metric is exploitable]{\textbf{Under the standard metric, less information
yields a higher score.} (a) A predictor that places every gene at the middle rank --- carrying
nothing about drug, cell or biology --- scores $\DEdr = 0.9999$, above the two-real-cell noise
ceiling and above both the model and its own scramble. The inverted ordering is the signature of an
artefact rather than of skill, and its cause is control regression-to-the-mean. Tier 2, $n=400$
conditions. No intervals are drawn: the stored ones are normal approximations rather than the
cell-line-clustered bootstraps used everywhere else in this thesis, and mixing the two conventions
on one axis would misrepresent both. (b) Grading the \emph{metrics} rather than the models. The
displacement drawn is the Dynamic Range Fraction, on a scale where $0$ is the drug-agnostic baseline
and $1$ a perfect predictor, so a negative value means the metric ranks an uninformative baseline
\emph{above} the noise ceiling. Only $\NIR$ points forward; the hardened estimate is
\ci{+0.635}{+0.618}{+0.652}. $655$ drugs over $61$ cell-line$\times$plate
groups, within plate. \texttt{weighted\_r2} carries no directional claim --- its sign tracks cell
count, and it is approximately neutral.}
\label{fig:calibration}
\end{figure}

The metric can discriminate even if the model cannot: a spike-in forced choice between two real
drug populations separates them essentially perfectly at pseudobulk ($\paneltau = 1.000$, $\DEdr =
0.995$, top-$n$ $\tau = 0.987$, plate held constant). The instrument is not the bottleneck.

\subsection{The model is drug-blind}
\label{sec:res-blind}

Having established which instrument can be trusted, we can finally ask the question the thesis
exists to answer. The answer is uniform across four instruments that fail in different ways, which
is the only reason to report all four.

\begin{table}[h]
\centering
\begin{tabular}{lrl}
\toprule
Instrument & Value & Reading \\
\midrule
Forced-choice grading & $0.48$ ($\approx$ scramble; ceiling 0.67--0.83) & chance \\
Output invariance (gap) & \ci{0.000}{-0.019}{+0.016} & null \\
Scramble $\DEdr$, tier 1 & $0.740$ real vs $0.739$ scrambled & identical \\
Scramble $\DEdr$, tier 2 & $0.733$ real vs $0.729$ scrambled & identical \\
$\gap$ (identifiable subset) & \ci{+0.014}{-0.016}{+0.042} & null \\
\bottomrule
\end{tabular}
\caption{Four independent instruments agree. Near-ceiling $\DEdr$ is achieved \emph{identically}
with a wrong-mechanism drug, so the score is drug-independent.}
\label{tab:blind}
\end{table}

Before asking what the model knows, it is worth asking what there is to know. A ladder of
mean-shift baselines --- each predicting the average response of a progressively narrower group ---
answers that directly. On tier 2, predicting the global mean scores $\DEdr$ $0.056$; conditioning
on cell line roughly triples it to $0.142$; conditioning on mechanism does not move it ($0.052$);
and conditioning on mechanism \emph{and} cell line is no better than cell line alone ($0.120$
against $0.142$). Tiers 1 and 3 tell the same story. The informative conditioning variable in this
data is therefore the cellular context, not the compound class --- worth holding onto, because
Section~\ref{sec:res-mechanism} finds precisely that asymmetry inside the model's activations, and
Section~\ref{sec:res-scope} finds that mechanism does carry signal once the generic program is
subtracted. The apparent tension is a frame difference rather than a contradiction: in
full-profile space the generic response dominates and swamps a mechanism effect that becomes
visible only in the residual frame.

On the calibrated metric the model sits at chance and is matched by predictors that know nothing
about the drug: within plate on tier 2 ($n = 606$), the replicate ceiling is $0.576$, the model
scores $0.498$, a drug-agnostic linear map $0.500$, a zero-drug-information control-copy $0.504$,
and the global mean $0.180$.

Two objections to this verdict close on measurement rather than argument. First, \emph{is the task
winnable at all?} Only $46.2\%$ of (drug, cell line) conditions are identifiable in principle;
$27.3\%$ of drugs are statistically inert against DMSO and $78.7\%$ have a plate-mate closer than
their own replicate. But restricting to the identifiable subset does not rescue the model: it
appears to beat the linear baseline there ($0.768$ vs $0.531$), yet a zero-drug-information
control-copy matches it exactly ($0.766$) and the leak-immune $\gap$ remains null. The model is
drug-blind \emph{even where the task is provably winnable}. The ceiling itself rises with
aggregation ($0.61$ at 10 cells, $0.76$ at 40, $0.85$ at 120), so the signal is real and
noise-limited.

\begin{figure}[htbp]
\centering
\includegraphics[width=\textwidth]{fig-difficulty}
\caption[Difficulty structure and the stratum ladder]{\textbf{The task is winnable where the noise
ceiling is high --- and precisely there, a predictor containing no drug information matches the
model.} (a) The ceiling rises steadily with the number of cells per condition, cross-plate and
within plate, so identifiability is noise-limited rather than absent: the drug signal is real and
emerges under aggregation. (b) Conditions binned by their own ceiling, with five arms in every
stratum. On the identifiable stratum ($n=204$) the model reaches $0.768$ --- and a control-copy
baseline that simply reproduces the untreated cell, and therefore knows nothing whatever about the
drug, reaches $0.766$. The paired contrast beneath each stratum is
\ci{+0.002}{-0.026}{+0.028} there, clustered over cell lines. Showing the model alone across these
strata would make the visual claim that it does better on winnable drugs; the control arm is what
shows that it does not.}
\label{fig:difficulty}
\end{figure}

Second, \emph{does the representation throw the drug away?} Cell sentences keep only the rank order
of genes and discard expression magnitude, and if the drug lived in the magnitudes the diagnosis of
the following sections would be attacking the wrong thing. The objection is testable without any
model, by asking whether \emph{true} expression discriminates drugs any better than rank does. It
does not: at single-cell resolution, cosine, Pearson and Spearman similarities on true normalised
expression all sit within $0.04$ of chance ($0.529$--$0.531$ [$0.515$, $0.543$]) --- the same
near-chance regime as rank. What recovers the signal is aggregation, not representation: at
fifteen cells the control-referenced shift cosine reaches $0.793$ [$0.775$, $0.809$]. The drug is
not hiding in the magnitudes, which matters twice over. It rules out the simplest alternative
explanation for everything that follows, and it means the repair of
Section~\ref{sec:res-encoding} cannot be read as ``restoring the magnitudes'': that repair keeps
rank throughout and changes only \emph{what is being ranked}. The same measurement bounds the task
from below --- at single-cell resolution the same-drug versus different-drug gap is $+0.002$, so
whatever a model can learn here, it cannot be learned from a single cell.

Third, \emph{does the scramble swap in similar drugs?} The swap-distance sweep resolves this:
$\gap$ is null in every distance bin including the far tail ($-0.019$), across a $2.3\times$ spread
in swap distance, with a trend correlation of $+0.051$. Telling the model it is a drug whose true
response is maximally different still does not change the output.

\subsection{Where the blindness lives}
\label{sec:res-mechanism}

A null result invites two readings, and they call for opposite responses: either the model cannot do
the task, or the task cannot be done. The previous section closed the second --- the ceiling rises
with aggregation, a third of conditions are identifiable in principle, and the model is at chance
even on those. So the failure is the model's, and the useful question becomes \emph{where} in the
model it happens. That question cannot be answered behaviourally; it needs the weights.

Drug identity \emph{is} in the representation: linear probes recover it at $82\%$ (layer 9) and
$76\%$ (layer 16) against a chance rate of $\approx 8\%$. With cell line, plate and dose projected
out, decodability \emph{rises} with depth, from $0.354$ at layer 2 to $0.883$ at layer 16. Yet the
purified drug subspace's variance share stays flat at $\approx 0.03$, against $\approx 0.6$ for cell
line.

\begin{table}[h]
\centering
\begin{tabular}{lrrrrr}
\toprule
Layer & Decode & Decode & Ablate pure & Random & Ratio \\
      & (raw)  & (context removed) & drug (KL) & matched dim & \\
\midrule
2  & 0.408 & 0.354 & 0.155 & 0.012 & $12.7\times$ \\
4  & 0.585 & 0.429 & 1.563 & 0.084 & $18.6\times$ \\
6  & 0.692 & 0.529 & 1.435 & 0.251 & $5.7\times$  \\
8  & 0.779 & 0.569 & 0.596 & 0.042 & $14.1\times$ \\
9  & \textbf{0.821} & 0.602 & 0.418 & 0.100 & $4.2\times$  \\
12 & 0.754 & 0.873 & 0.133 & 0.035 & $3.8\times$  \\
16 & 0.758 & \textbf{0.883} & 0.015 & 0.004 & $3.6\times$  \\
\bottomrule
\end{tabular}
\caption{All seven layers pass the removal gate (88--95\% of above-chance drug signal removed);
chance is $0.083$. Two decode series are shown because they behave differently and the distinction
matters: the \emph{raw} probe peaks at layer 9 and then declines, whereas with cell line, plate and
dose projected out decodability rises monotonically to layer 16. The identity test is reported
separately below, on a later and differently controlled design; the ablation ratios here are a
statement about \emph{functional variance} and carry no claim that drug identity is read.}
\label{tab:workspace}
\end{table}

\begin{figure}[htbp]
\centering
\includegraphics[width=\textwidth]{fig-mechanism}
\caption[The drug is encoded more with depth and never read for identity]{\textbf{The model computes
an increasingly explicit drug representation and sets it aside.} (a) With cell line, plate and dose
projected out, drug decodability rises monotonically with depth, from $0.354$ at layer 2 to $0.883$
at layer 16, against a chance rate of $0.083$; the causal variance share of that same purified
subspace stays flat near $0.03$, against roughly $0.6$ for cell line. The raw probe is drawn too and
behaves differently --- it peaks at layer 9 and declines --- because the two series are easily
conflated and the dissociation rests on the context-removed one. (b) Ablation of the purified drug
subspace against a matched-\emph{dimension} random subspace, on a log axis because late-layer
divergences are small for everything. This is a functional-variance comparison: it establishes that
the subspace is not inert, and it makes no claim about whether drug identity is read, which requires
the substitution test reported in the text.}
\label{fig:mechanism}
\end{figure}

The identity test in this section has been rebuilt twice, and the failures are reported because
they bear on how much the replacement can carry. The first compared a drug-B injection inside a
$\leq 10$-dimensional purified subspace against a random vector drawn in the full
$2048$-dimensional stream, so roughly $99.5\%$ of the control's energy lay in directions the
intervention never touched. The second drew the control inside the subspace, which looked like a
repair and was not: the hook ablates before adding, so what is delivered is (added $-$ removed),
and an uncentred class mean shares its global-mean component with the thing removed. The
substitution arm was then close to a restoration and the control arm a real displacement. Simulated
with no model present, the control delivers $5.3\times$ the perturbation in $100\%$ of draws, with
the predicted gap $2\langle c_B, s\rangle = 9.179$ against an observed $9.162$. Both retracted
results said the substitution moves the output \emph{less} than the control, and both were
reporting geometry. What survives untouched is everything not resting on the substitution: the
ablation ratios, whose null was always a matched-\emph{dimension} random subspace and correctly
constructed; the variance shares; and the decodability-versus-depth dissociation in
Table~\ref{tab:workspace}.

The replacement design of Section~\ref{sec:methods-probe} was run across five training arms and,
for the two residual arms, replicated at three further seeds. Of $26$ headline tests --- five arms
by up to seven depths --- exactly one survives Benjamini--Hochberg correction: the residual arm at
layer $12$.

\begin{table}[h]
\centering
\begin{tabular}{lrrr}
\toprule
Arm / seed & effect & 95\% CI & $p$ \\
\midrule
residual, seed 42 & $+0.0197$ & $[+0.0063, +0.0335]$ & $0.0010$ \\
residual, seed 43 & $+0.0118$ & $[+0.0060, +0.0173]$ & $0.0003$ \\
residual, seed 45 & $+0.0085$ & $[+0.0029, +0.0143]$ & $0.0015$ \\
\midrule
single-cell arm, six depths & \multicolumn{3}{l}{nothing detectable at any depth (median $-0.0012$)} \\
\bottomrule
\end{tabular}
\caption{The identity test, as a fraction of that model's own clean preference gap between the two
responses being compared. Seed 44 produced no layer-12 measurement --- the instrument was excluded
there for lack of dynamic range --- so the pre-registered criterion of two replications in three new
seeds is met with the third a non-measurement rather than a failure.}
\label{tab:identity}
\end{table}

The value of this is that it is a \emph{positive control on real data}. The residual arm is
independently known to use the drug: it is the arm whose prompt sensitivity is established
behaviourally in Section~\ref{sec:res-generalisation}. The probe detects identity reading there,
replicably, and detects none anywhere in the arm that is behaviourally drug-blind. A null from an
instrument that has never been shown to register a positive is uninterpretable; this one has.

The claim the evidence supports is that the generation readout consults the drug representation
\emph{at very low gain}: about $1\%$ of the preference gap, an order of magnitude below the
pre-registered margin for calling an effect material, in the one arm where it is detectable at all.
It is not that the readout cannot read drug identity, and it is not the stronger mechanistic story
that earlier drafts of this thesis told. Three limits belong with it. The effect appears at one
depth of seven, and layer 4 is positive in a single seed and null in the other three, which is what
multiplicity discipline is for. The single-cell arm has no layer-12 measurement at all --- it was
excluded there for instrument saturation --- so the cleanest comparison, the same depth across two
arms, does not exist; the dissociation rests on that arm being null at all six depths it could
measure. And the sibling residual-holdout arm does not reproduce the effect. Read together with
the decodability series, the picture is a representation that becomes more explicit with depth
while the process that writes the output barely consults it --- which is suggestive of the
workspace account of Section~\ref{sec:lit-interp}, and is offered as an interpretation of the
pattern rather than as a demonstration of the mechanism.

\subsection{The $\varepsilon$-ladder: no target-side fix works}
\label{sec:res-epsilon}

If the defect is a readout that feels the magnitude of activity in the drug region but not its
direction, then improving what sits \emph{in} that region cannot help --- and improving the training
target is exactly that kind of intervention. The obvious objection is that one failed target
construction proves nothing, so rather than try another we swept the family: entropic
regularisation gives a one-parameter path whose two endpoints are the two constructions already
tested independently, which turns a sequence of attempts into a controlled sweep.

\begin{table}[h]
\centering
\begin{tabular}{llrl}
\toprule
Target & $\varepsilon$ & $\gap$ & Verdict \\
\midrule
single cell (real treated cell) & $\approx 0$ & \ci{+0.014}{-0.016}{+0.042} & null \\
OT barycentric (state-matched) & $0.5$ & \ci{-0.001}{-0.018}{+0.016} & null \\
consensus (global mean) & $\rightarrow\infty$ & \ci{-0.010}{-0.022}{+0.003} & null \\
\bottomrule
\end{tabular}
\caption{The regularisation parameter interpolates between the two target constructions tested
independently, and no point on the ladder lifts the model out of the near-null regime. These figures
come from a \emph{random-partner} scramble; Section~\ref{sec:res-encoding} shows that under a
sharper stratified test the optimal-transport arm does show a weak effect
(\ci{+0.0263}{+0.0069}{+0.0476}), roughly five-fold below what re-encoding the target achieves.
The consensus endpoint is additionally the ladder's weakest-evidenced point: that arm was stopped at
roughly $60\%$ of an epoch and its tier-1 NIR was unscoreable, so the negative rests on the
single-cell and optimal-transport endpoints.}
\label{tab:epsilon}
\end{table}

Training was not inert: the consensus model stops echoing the control ($0.653$ against a
control-copy of $0.766$, where the single-cell model sits \emph{on} the control at $0.771$) and its
predictions are not mode-collapsed. But the variation is uncorrelated with real drug identity
(Mantel $\approx +0.03$). Not collapse --- wrong structure. Denoising changes \emph{what is in} the
target while the defect is what its tokens encode; that is the quantity the next two sections
measure and then change.

\subsection{How much signal is there to encode? The transfer coefficient}
\label{sec:res-transfer}

Before re-encoding anything, it is worth knowing how much drug-specific signal exists and what kind
it is, because the answer scopes every claim in this thesis. The residual decomposes into a
per-drug constant $\beta(d)$ and a drug$\times$cell-line interaction $\kappa(d,c)$
(Section~\ref{sec:methods-vardecomp}), and the two components call for different responses: the
first is retrievable without any model, the second is the only thing a conditional model could win
on. $\NIR$ cannot adjudicate between them, so the question is posed in cosine space with the
measurement noise divided out.

The measurement itself carries a control-design lesson, stated because it changed the protocol.
The first run compared $\Tcoef$ against a negative control of \emph{different drugs on the same
plate}, which returned $+0.485$ against $\Tcoef = 0.511$, and was correctly declared void: that
control differs from the estimand in two ways at once --- different drug \emph{and} same cell line
--- and plate structure surviving a cell-line-scoped generic inflates same-plate comparisons (to
$+0.478$) while leaving cross-line comparisons untouched. A \emph{structure-matched} control ---
different drug, different cell line, breaking only the drug match --- is the one that bounds the
estimand, and a negative control that differs in more than one way cannot do so.

\begin{table}[h]
\centering
\begin{tabular}{lrrr}
\toprule
Quantity & Estimate & 95\% interval & Reference \\
\midrule
$\Tcoef$ (cross cell lines, plate-scoped generic) & $0.558$ & $[0.514, 0.593]$ & simulated $\kappa{=}0$ null $\approx 1.0$ \\
drug $\times$ cell-line interaction share & $44.2\%$ & $[40.7\%, 48.6\%]$ & --- \\
$\Tcoef$ across dose (same drug, same line) & $0.707$ & $[0.632, 0.782]$ & reference scale \\
matched negative controls & $-0.017$, $-0.008$ & spanning zero & must be $\approx 0$ \\
\bottomrule
\end{tabular}
\caption{The transfer coefficient, estimated on real molar doses with dyadic cluster-robust
intervals (Section~\ref{sec:methods-vardecomp}). The matched negative control sits on zero; the
same-plate control that produced the earlier false alarm reads $+0.478$ for the reasons given in
the text.}
\label{tab:transfer}
\end{table}

Roughly $44\%$ of the drug-specific residual variance is drug$\times$cell-line interaction, robust
to the generic's scope and to control splitting. Two internal checks agree, and both were decisive
between scopes. \emph{Dose ordering}: same-drug, same-line, different-dose transfer is $0.707$
[$0.632$, $0.782$], above the cross-line $0.558$ --- changing the dose costs \emph{less} than
changing the cell line, the biologically correct ordering; at cell-line scope the ordering
inverts. \emph{Control-split bias}: at plate scope the shared- and split-control builds agree to
$-0.000$, against $+0.254$ at cell-line scope, so that discrepancy was plate contamination rather
than the split. The reliability that enters the disattenuation is a within-well sampling precision
(Section~\ref{sec:methods-residual}), optimistic relative to true biological reproducibility; the
optimism sits in the denominator, so the interaction share is if anything understated. The per-drug
constant $\beta(d)$ --- which the lookup of Section~\ref{sec:res-lookup} reproduces almost exactly
--- is therefore \emph{not} the whole signal. There is a conditional target worth about half the
drug-specific variance that a per-drug constant cannot express.

Whether anything could \emph{learn} that interaction is a separate question, and the direct test is
sobering. Writing $\kappa(d,c) = \resid(d,c) - \hat{\beta}(d)$ with $\hat{\beta}$ estimated
leave-one-cell-line-out, and disattenuating by $\kappa$'s own split-half reliability:

\begin{table}[h]
\centering
\begin{tabular}{lrrr}
\toprule
Configuration & same cell line & different cell lines & excess within-line \\
\midrule
plate-scoped generic & \ci{-0.004}{-0.011}{+0.002} & $+0.003$ & $-0.007$ \\
cell-line-scoped & \ci{-0.003}{-0.021}{+0.015} & $+0.002$ & $-0.005$ \\
\bottomrule
\end{tabular}
\caption{No cell line bends different drugs in a consistent direction. $4{,}000$ pairs per arm; the
signal interval at plate scope excludes anything above $+0.002$, so this is a real null rather than
an underpowered one. The cell-line-scoped row is wider ($[-0.021, +0.015]$) and carries the claim
less strongly.}
\label{tab:kappa}
\end{table}

We looked for structure in the interaction and did not find it: no cell line modulates different
drugs in a consistent direction, and a direct test of whether mechanism-matched drugs share an
interaction with a given cell line returned nothing detectable (permutation $p \approx
0.15$--$0.31$). The temptation is to read two negative results as \emph{irreducibility}, and they
do not support it. The second test bounds its own power severely: $\kappa$'s split-half reliability
in this atlas is $0.195$, so the quantity being correlated is barely estimable, and same-mechanism
compounds are so heavily co-plated that only $3$--$7\%$ of matched pairs survive a different-plate
requirement, which leaves the plate confound in place. The defensible reading is that $\kappa$'s
structure is \emph{unresolved in this dataset}. Resolving it would take more cells per condition,
or a plating design that crosses mechanism with plate --- neither of which is an analysis choice.

The same measurement scopes the whole thesis. The shift's energy decomposes exactly into residual
$0.621$, generic $0.390$ and a cross term of $-0.011$: $62.1\%$ of the perturbation response's
energy is drug-specific, $39.0\%$ is the generic program, and the two components are essentially
orthogonal, so the shares are genuine. Every claim in this thesis is scoped by the first of those
numbers, and it has not previously been computed for this task. It also forces a retraction: the
justification for estimating the generic per cell line rather than per plate was a reproducibility
comparison, $62\%$ against $19\%$, measured with \emph{shared} control cells --- which inflate
cell-line reproducibility far more than plate reproducibility. Measured with split controls the two
are comparable ($20\%$ plate, $16\%$ cell line), so the argument favours neither, and the
dose-ordering and control-split diagnostics above favour plate outright. The production targets are
therefore built at plate scope (Section~\ref{sec:methods-residual}).

\subsection{Changing what the tokens encode}
\label{sec:res-encoding}

Every target-side intervention changed \emph{what the target was} and none changed \emph{what its
tokens encode}. That distinction is measurable rather than rhetorical: because a cell sentence
ranks genes by expression, the target strings for two drugs in the same context are nearly
identical, and the loss weights every token position equally.

\begin{table}[h]
\centering
\begin{tabular}{lrr}
\toprule
Representation & Inter-drug top-200 tokens differing & Replicate retrieval ceiling \\
\midrule
full profile (the standard target) & \textbf{34.6 / 200 (17\%)} & 0.742 \\
shift (treated $-$ control) & 111.4 (56\%) & 0.742 \\
\textbf{residual (drug-specific)} & \textbf{118.2 / 200 (59\%)} & 0.7432 \\
\bottomrule
\end{tabular}
\caption{$83\%$ of the standard target's tokens are identical between any two drugs in the same
context. The residual row is quoted at \textbf{cell-line} scope; at plate scope it reads $120.3$
and $0.7467$. Note the retrieval ceiling is \emph{unchanged} across all three rows: the information
was always present, and the tokenisation discarded it.}
\label{tab:divergence}
\end{table}

The natural inference is that only a small fraction of the gradient can be carrying drug identity.
That inference is a \emph{hypothesis} about the optimisation: what is measured here is the overlap
of the target gene sets, not a token-level loss or gradient attribution, and it should be read at
that strength. Its warrant is predictive --- it explains why every change to target \emph{content}
failed identically, and it predicts that changing what the tokens \emph{encode} would not.
Re-encoding the target as the drug-specific residual raises the differing fraction to $59\%$ at an
unchanged retrieval ceiling, and produces the first non-null result in the project. On the earlier
target build, the stratified scramble gave a gap that grew with swap dissimilarity ---
\ci{+0.0351}{+0.0061}{+0.0656}, \ci{+0.0716}{+0.0371}{+0.1059} and \ci{+0.1429}{+0.1112}{+0.1787}
across the \texttt{near}, \texttt{orth} and \texttt{opposite} strata --- while the single-cell and
optimal-transport controls stayed flat and null at every stratum ($-0.001$ to $+0.016$). The
monotone gradient is itself the evidence: a metric artefact would not order itself by swap
dissimilarity.

\begin{figure}[htbp]
\centering
\includegraphics[width=\textwidth]{fig-repair}
\caption[The gap grows with swap dissimilarity]{\textbf{On the earlier target build, the residual
model's advantage over its own scramble grew with how dissimilar the swapped-in drug's response was;
a control scored through the identical pipeline stayed flat and null.} The evidence is the gradient,
not any single point --- a metric artefact would not order itself by swap dissimilarity. The
horizontal axis is the measured mean residual cosine between the real and swapped-in drug, which is
byte-identical across all four runs, so the positions are exact. The three solid series are at a
matched $600$-token generation budget; the residual model's $1400$-token run is drawn separately as
the dotted overlay, since the budget alone doubles the effect. Intervals are $95\%$ bootstrap
clustered over cell lines, $n=250$ conditions per point. The \texttt{near} stratum is the weakest
test by construction: where the swapped-in drug's true response is similar, an unchanged output is
correct rather than blind. These magnitudes are superseded by the rebuilt-target run of
Section~\ref{sec:res-generalisation}; the figure stands as the first demonstration of the effect.}
\label{fig:repair}
\end{figure}

Three further diagnostics on that build agreed, and only the residual model passed all three: the
correlation between condition reproducibility and model $\NIR$ was $+0.144$ (the model does better
where the drug effect is real) against $-0.117$ and $-0.216$ for the controls; prediction diversity
was highest for the residual model; and while $\cos(\textrm{prediction}, \textrm{own truth})$
exceeded $\cos(\textrm{prediction}, \textrm{other drugs})$ for all three arms, the residual model's
margin ($+0.084$) was roughly twice the optimal-transport arm's ($+0.049$) and three times the
single-cell arm's ($+0.026$).

Two measurement lessons belong with this result. The effect was first measured at a $600$-token
generation cap, which silently truncated $26\%$ of generations --- a signature is $\approx 201$ gene
symbols at 2--4 BPE tokens each --- and \emph{halved} it ($+0.0716$ against the correct $+0.1429$);
re-running both controls at the matched $1400$-token budget left them flat and null, so the
comparison is budget-matched and the conclusion unaffected. And the optimal-transport arm,
originally reported as achieving nothing, scores \ci{+0.0263}{+0.0069}{+0.0476} under the sharper
stratified test with a monotone gradient: it does use the drug, weakly. The original null was an
underpowered instrument, not an absence of signal; the $\varepsilon$-ladder conclusion stands, but
``OT achieved nothing'' is withdrawn.

What those magnitudes can and cannot claim is re-adjudicated in the next section, because they were
read against the \texttt{opposite} comparator, and that comparator turns out not to be a null. The
comparator-free evidence is already in hand, and it is the strongest in this thesis: regressing each
generation's $\NIR$ on the cosine between the named drug's truth and the target's truth gives a
slope of \ci{+0.390}{+0.225}{+0.542} on trained conditions --- zero slope is the natural null, so
no comparator defect can reach it. Re-encoding the target made the model respond to the combined
drug-and-mechanism prompt. What that response is worth is the question the rest of the chapter
answers.

\subsection{Sensitivity is not fidelity: a main effect, not an interaction}
\label{sec:res-generalisation}

A model that responds to the drug token has cleared a low bar, because reproducing a training
average requires responding to it too. Everything scored above came from conditions the model
trained on, so the result so far shows that the token is consulted --- not that anything
transferable was learned. Separating those two requires withholding data at the right granularity
(Section~\ref{sec:methods-holdout}), and it requires a comparator that is genuinely uninformative.
The first finding of this section is that the comparator used up to here was not, and every
magnitude quoted against it must be re-read.

The comparator's own score gives it away. Across the three swap strata, the scrambled model's
$\NIR$ is $0.572$, $0.497$ and $0.403$ as the named partner's true response moves from aligned
($\cos = +0.26$) through orthogonal ($0.00$) to anti-aligned ($-0.24$) with the target. A
comparator whose score is a monotone function of which drug it was handed is not a null: it is an
active treatment that tracks the partner drug. The swap partner is itself a drug the model was
mostly trained on, so telling the model a lie costs it precisely because it knows the partner. A
gap measured against such an arm decomposes as
\begin{equation}
  \gap \;=\; \underbrace{\big(\NIR_{\textrm{model}} - 0.5\big)}_{\textrm{how much naming the truth helps}}
  \;+\; \underbrace{\big(0.5 - \NIR_{\textrm{scramble}}\big)}_{\textrm{how much naming the lie hurts}},
\end{equation}
and only the first term is wanted. Against \texttt{opposite} the second term is $+0.097$; against
\texttt{orth}, whose own score sits on chance ($0.497$), it is zero. Only the neutral column of
Table~\ref{tab:generalisation} can therefore carry a magnitude claim, and the
\texttt{opposite} column is reported solely as the diagnostic that would have misled us.

\begin{table}[h]
\centering
\begin{tabular}{lrrll}
\toprule
Split & $n$ (cell lines) & model $\NIR$ & $\gap$ vs \texttt{orth} (neutral) & $\gap$ vs \texttt{opposite} \\
\midrule
\texttt{train} & 200 (35) & 0.635 [0.575, 0.695] & \ci{+0.1457}{+0.0836}{+0.2078} & \ci{+0.2627}{+0.2100}{+0.3155} \\
\textbf{\texttt{unseen\_combo}} & \textbf{250 (39)} & \textbf{0.545 [0.493, 0.596]} & \ci{+0.0449}{-0.0074}{+0.0972} & \ci{+0.1281}{+0.0768}{+0.1795} \\
\texttt{unseen\_drug} & 120 (38) & 0.512 [0.442, 0.581] & \ci{+0.0097}{-0.0655}{+0.0848} & \ci{+0.0862}{-0.0028}{+0.1753} \\
\bottomrule
\end{tabular}
\caption{Whole-well holdout on the rebuilt targets (split-before-fit, plate-scoped generic). Only
the \texttt{orth} column can carry a magnitude claim; the \texttt{opposite} column adds ``how much
naming a lie hurts'' to ``how much naming the truth helps'' and is shown as the diagnostic. The
\texttt{unseen\_drug} row is the control that proves the point: those are drugs the model has never
seen, its gap must be zero, and it is zero only in the neutral column. Model $\NIR$ intervals test
against the chance value $0.5$ with no comparator. Intervals are two-way cluster-robust over cell
lines and wells, with a Student $t$ reference (Section~\ref{sec:methods-controls}).}
\label{tab:generalisation}
\end{table}

The control behaves exactly as the diagnosis predicts. On \texttt{unseen\_drug} --- drugs held out
of fine-tuning, where no learned knowledge can be at work --- the gap against \texttt{opposite} is
$+0.0862$ \nolinebreak[3]$[-0.0028, +0.1753]$, nine times the neutral estimate, with an interval
that only just contains zero. The same conditions, the same generations and the same model scored
against the neutral stratum give $+0.0097$ \nolinebreak[3]$[-0.0655, +0.0848]$, sitting on zero as
a control should. The borderline interval on the \texttt{opposite} arm is itself a lesson in the
inference regime: a cluster-robust variance is asymptotic in the number of \emph{clusters} rather
than observations, this arm spans only $25$ treatment wells, and under the correct $t_{24}$
reference ($2.06$ rather than $1.96$) the interval widens by five per cent and covers zero. An
earlier draft of this section read the normal-quantile interval as excluding zero; it does not.

Read in the neutral column, \texttt{train} clears zero at \ci{+0.1457}{+0.0836}{+0.2078} and
\texttt{unseen\_combo} does not: $+0.0449$ \nolinebreak[3]$[-0.0074, +0.0972]$, spanning zero
under well-clustering and line-clustering alike ($[-0.0050, +0.0993]$). Generalisation to an unseen
treatment well of a seen drug is therefore \emph{not established}. The precision matters: not
established is not the same as shown to be zero. The interval reaches $+0.097$, so an effect of up
to roughly $+0.1$ remains compatible with these data at $n = 250$ --- a limit of power at the
sample size the holdout affords, reported as one. (On \texttt{unseen\_drug} the corresponding
detection limit is $+0.099$.) The difficulty of the three splits is not the explanation: their
split-half replicate ceilings are $0.962$, $0.955$ and $0.956$, so the splits are equally hard and
the held-out shortfall is not a noise-floor difference.

The difference between the trained and held-out estimates is $+0.1008$
\nolinebreak[3]$[+0.0196, +0.1820]$ with cluster-robust $p = 0.0150$ (an unclustered permutation
returns $p = 0.0054$, quoted only for contrast, because ignoring the clustering this thesis
mandates everywhere else flatters the result).
% AUDITOR NOTE: the [+0.0196, +0.1820] interval is carried over from the pre-revision text; the
% fixed-claims register fixes only the point estimate and p. If the interval was computed under a
% normal rather than Student-t reference it should be recomputed before submission. What that
advantage \emph{is} requires care, because the two arms differ in two ways at once: whether the
condition was trained on, and \emph{which drugs they contain}. Restricting the comparison to the
$20$ drugs present in both arms, the advantage attenuates to $+0.0420$ with $p = 0.4764$. The
defensible statement is therefore narrow: a condition-weighted training-exposure advantage exists
in this evaluation, and its generality across drugs is unresolved --- the design cannot separate
exposure from drug identity, and we do not claim either direction of the stronger interpretation.

Two instruments that need no comparator at all cannot be touched by any of this, and they carry the
positive content of the section. First, the slope test: regressing each generation's $\NIR$ on the
cosine between the named drug's truth and the target's truth gives \ci{+0.390}{+0.225}{+0.542} on
\texttt{train}, \ci{+0.374}{+0.214}{+0.508} on \texttt{unseen\_combo} and
\ci{+0.238}{+0.033}{+0.420} on \texttt{unseen\_drug} --- all three exclude zero. The model emits
something aligned with the named drug's signature on every split, including held-out wells, with no
degradation worth the name between train and held-out. Second, scored against chance with no
comparator, the model's own $\NIR$ clears $0.5$ on \texttt{train} ($0.635$ [$0.575$, $0.695$]) and
is not distinguishable from it on either held-out split. A positive slope alongside a chance-level
$\NIR$ is not a contradiction: it says the drug-to-signature map exists but is \emph{coarse}. The
model knows roughly what a drug does and not what it does \emph{in this well} --- and $\NIR$ asks
the second question, ranking against the $\approx 40$ other drugs measured in the same cell line.
That is the drug main effect without the drug $\times$ cell-line interaction, which
Section~\ref{sec:res-transfer} measures at $44.2\%$ of the residual variance. It is also why the
lookup wins in the next section: \texttt{drug\_lookup} \emph{is} that main effect, estimated
straight from data far more precisely than the model learned it.

The directional-agreement measurement says the same thing at condition level. The cosine between
the model's prediction and its own condition's truth on trained conditions averages $+0.0511$
(sd $0.1141$, median $+0.0300$), with individual conditions spanning $[-0.1819, +0.4395]$ and
$32.0\%$ clearing $|0.1|$: weak on average and substantially heterogeneous, so a single summary
number would misrepresent it in either direction.

One caveat attaches to every held-out number in this chapter, and it is stated here where those
numbers live. The evaluation's truths are reconstructed through a compatibility path that refits
the generic on the reliability-surviving subset of training conditions rather than the full set the
checkpoint was trained against, and that filters held-out conditions by their own split-half
reproducibility --- a selection on the evaluation set's outcome. Its direction is conservative for
the headline: it retains the cleanest, easiest conditions, and the held-out result still fails to
clear zero. But the evaluated population is not the full held-out population, and the intervals
above should be read as applying to the reproducible subset.

Two instrument defects, both of which would have concealed everything in this section, are worth
recording because they are now checks rather than assumptions. A uniform condition sampler
reproduces the pool's split proportions, so $250$ constructed held-out combinations yielded only
$33$ scored ones; and a strict-inequality $\NIR$ scored the zero vector (``predict the generic
response'') at $0.000$ rather than $0.500$. Both are corrected (Section~\ref{sec:methods}), and
after correction the drug-agnostic baselines land where they must: \texttt{generic} exactly
$0.500$, \texttt{control\_copy} $0.511$, \texttt{random} $0.521$. The residual frame is
\emph{demonstrably} leak-proof rather than assumed to be.

The sentence the evidence supports, then, is narrower than the one this section originally carried:
re-encoding the target made the model respond to the combined drug-and-mechanism prompt and learn a
drug-level average response. It learned the main effect and not the interaction, and its
generalisation to unseen wells is a power-limited open question rather than a result in either
direction.

\subsection{But a lookup table wins}
\label{sec:res-lookup}

Transfer without fidelity is a real property, but it says nothing about quality: a model can apply
a learned signature in a new context and still apply it badly. The pre-registered bar was therefore
never ``beats chance'' --- residual targets are plausibly memorisable, and a predictor that simply
looks up what this drug did elsewhere needs no model at all. That predictor is the comparison that
matters.

\begin{table}[h]
\centering
\begin{tabular}{lrl}
\toprule
Arm & $\NIR$ & note \\
\midrule
ceiling (within-condition split-half replicate) & $0.958$ & $0.962$ / $0.955$ / $0.956$ by split \\
\texttt{drug\_lookup} (training-only) & $0.941$ & the bar to beat \\
\texttt{drug\_lookup\_1} (one other line) & $0.859$ & removes the averaging advantage \\
model & $0.569$ & $0.635$ / $0.545$ / $0.512$ by split \\
\texttt{random} & $0.521$ & floor \\
\texttt{control\_copy} & $0.511$ & drug-agnostic \\
\texttt{generic} & $0.500$ & the zero vector, by construction \\
\bottomrule
\end{tabular}
\caption{The baseline ladder on the rebuilt targets, $n = 570$ scored conditions; the lookups are
fitted from training conditions only (Section~\ref{sec:methods-baselines}). The paired comparison
on the common support ($n = 449$) is model $-$ \texttt{drug\_lookup} $= \ci{-0.3559}{-0.3988}{-0.3130}$,
and of the distance from chance to the replicate ceiling the model covers $18.6\%$ against the
lookup's $96.3\%$.}
\label{tab:lookup}
\end{table}

The pre-registered interpretation, written before the result, anticipated three branches; the middle
one fired. The model reads and transfers the drug, but recovers only a small fraction of its
signature. On the earlier build it also beat \texttt{moa\_lookup} ($0.529$, $n = 143$), so its
prompt response is not reducible to the \texttt{Mechanism:} field --- though at that level the bar
is low.

This is DrEval's finding, in transcriptomic space. DrEval reports that deep learning models barely
outperform a naive predictor of mean drug and cell-line effects~\cite{dreval}. Their
\texttt{NaiveMeanEffectsPredictor} is $\hat{y}_{ij} = \mu^c_i + \mu^d_j - \mu$; our
\texttt{drug\_lookup} is the $\beta(d)$ term of that same additive model, measured in a frame where
the $\textrm{generic}(c)$ term has already been subtracted. The finding replicates and sharpens: on
a $946$-gene transcriptomic response rather than a scalar potency, and against a proper
within-condition replicate ceiling, the additive drug main effect reaches $0.941$ of a $0.958$
ceiling while a $1$B-parameter model reaches $0.569$. Note also that the warning DrEval inherits
from~\cite{partin} --- that good pair-holdout performance is reachable by predicting each drug's
training average --- is treated here as the hypothesis rather than the caveat, and measured. It
holds.

It is tempting to read $\textrm{ceiling} - \texttt{drug\_lookup} = +0.017$ as ``the modelling
target has vanished''. That inference is not supported, for three reasons that all favour the
lookup: the ceiling is scored against a \emph{half-sample} truth while every baseline is scored
against the \emph{full-sample} truth; \texttt{drug\_lookup} averages $\approx 38$ conditions against
the ceiling's single noisy half; and $\NIR$ near $0.96$ is compressive, collapsing a wide range of
cosines into a few thousandths. The paired gap on common support, $-0.3559$
\nolinebreak[3]$[-0.3988, -0.3130]$, is the honest summary of the contest; the coverage figures,
$18.6\%$ against $96.3\%$ of the chance-to-ceiling distance on that same support, are its
scale-free restatement.

\subsection{Every arm in one place}
\label{sec:res-summary}

The results so far are distributed across two measurement frames, and a reader assembling them has
to keep track of which numbers may be compared with which. Table~\ref{tab:allarms} collects them.

The two frames are \emph{not} interchangeable and are separated deliberately. The expression frame
scores a predicted cell profile against other drugs on the same plate; the residual frame scores a
predicted drug-specific signature against other drugs in the same cell line, after the generic
program has been subtracted. Both have chance at $0.50$ by construction, and that coincidence is the
trap: the residual frame's ceiling is $0.958$ where the expression frame's is $0.576$, so a model
scoring $0.569$ in one and $0.498$ in the other has not improved --- it is being measured against a
different question.

\begin{table}[h]
\centering
\begin{tabular}{llrl}
\toprule
Frame & Arm & $\NIR$ & What it knows \\
\midrule
\multirow{5}{*}{\shortstack[l]{expression\\(within plate,\\tier 2, $n=606$)}}
 & noise ceiling & 0.576 & a real replicate of this condition \\
 & control-copy & 0.504 & nothing --- it reproduces the untreated cell \\
 & linear map & 0.500 & the control cell, no drug \\
 & \textbf{model} & \textbf{0.498} & control cell $+$ drug \\
 & global mean & 0.180 & nothing \\
\midrule
\multirow{7}{*}{\shortstack[l]{residual\\(cell-line sets,\\$n=570$)}}
 & noise ceiling & 0.958 & a real split-half replicate of this condition \\
 & drug lookup (training-only) & 0.941 & this drug's signature from other cell lines \\
 & drug lookup, one line & 0.859 & the same, from a single other cell line \\
 & \textbf{model} & \textbf{0.569} & control cell $+$ drug \\
 & scramble, \texttt{orth} partner & 0.497 & the neutral comparator: chance by construction \\
 & scramble, \texttt{opposite} partner & 0.403 & an active treatment, not a null \\
 & generic response & 0.500 & nothing (the zero vector, by construction) \\
\bottomrule
\end{tabular}
\caption{Every arm on the calibrated metric, by frame. Chance is $0.50$ in both. The residual-frame
figures are the rebuilt-target run of Sections~\ref{sec:res-generalisation}
and~\ref{sec:res-lookup}; the held-out subset carries the outcome-selection caveat stated there.}
\label{tab:allarms}
\end{table}

Read down the residual block and the thesis's shape is visible in one object. A replicate of the
experiment reaches $0.958$; a lookup table reaches $0.941$; the model reaches $0.569$; the same
model told a response-orthogonal drug sits at $0.497$, the neutral zero point, and told an
anti-aligned drug falls to $0.403$. The distance between the model and its neutral scramble is the
drug effect this work recovered; the distance between the model and the lookup is what it did not.

\subsection{Scope: what is closed by measurement}
\label{sec:res-scope}

An unseen drug can only be predicted through a property it \emph{shares} with drugs seen in
training, so the question is answerable by enumerating the channels and measuring each in the same
frame. For each channel we predict a drug's residual from the mean residual of the \emph{other}
drugs in the same cell line that share the property --- the construction of \texttt{moa\_lookup},
generalised --- and pair it with a plate-matched random null, because the estimand requires the two
arms to differ in exactly one respect.

\begin{table}[h]
\centering
\begin{tabular}{lrrrrl}
\toprule
Channel & Coverage & $n$ & vs plate-matched null & (vs count-matched) & Verdict \\
\midrule
\textbf{protein target} & 54\% & 894 & \ci{+0.0973}{+0.0741}{+0.1183} & $+0.1269$ & \textbf{live} \\
\textbf{mechanism (MoA)} & 32\% & 1053 & \ci{+0.0749}{+0.0538}{+0.0948} & $+0.1067$ & \textbf{live} \\
chemical structure & 100\% & 4084 & $+0.0162$ [${+0.0051}$, ${+0.0264}$] & $+0.0206$ & closed \\
\bottomrule
\end{tabular}
% AUDITOR NOTE: the coverage and n columns are carried over from the count-matched run; the
% plate-matched re-run's register does not restate them. They should be near-identical (same
% annotation support) but ought to be re-quoted from the corrected artifact before submission.
\caption{Coverage is stated first: a channel annotated on few drugs cannot be \emph{closed} by a
null result, because it was never tested. The plate-matched null is the estimand-correct control;
the count-matched column shows how much of each raw gap was co-plating. Chemistry reproduces the
structure gate in this frame; both knowledge channels clear their corrected nulls.}
\label{tab:channels}
\end{table}

The size of that correction is itself informative, so it is reported rather than argued away.
Mechanism partners are co-plated with the target $0.420$ of the time against $0.340$ for a random
draw (protein target $0.348$ against $0.323$; chemical similarity $0.344$ against $0.345$ --- none,
as expected for a channel uncorrelated with plating). Under a third construction that draws all
partners from \emph{other} plates, the knowledge channels' advantages fall to $+0.0643$ and
$+0.0344$ and chemistry's interval spans zero ($+0.0056$); and in that construction the channel
arms' absolute $\NIR$ sits at $0.360$ and $0.296$, \emph{below} chance --- a direct demonstration of
how strongly the cell-line-scoped residual frame is dominated by plate. The verdicts above survive
the correction; the magnitudes shrank by a fifth to a third, which is what the co-plating fractions
predict.

This overturns a pre-registration made in this project. Two design documents concluded that
unseen-drug generalisation was out of reach, and both inferred it from the structure gate --- that
is, from Morgan fingerprints and a molecular language model. But the drug-knowledge specification
written earlier in the same project ranks raw structure the \emph{lowest}-value channel and protein
targets ``the key feature'', on the reasoning that transcriptional response is driven by which
protein a compound hits, something structure encodes only indirectly. The pre-registration
generalised from the weakest channel to the strongest, and it was wrong. Tahoe's own metadata
carries protein-target annotations for $264$ drugs --- $280$ distinct symbols, $550$ edges ---
which no analysis in this project had read.

The correct statement is therefore narrower and more useful than the one we started with:
\emph{chemical structure does not predict response, measured twice in two frames; protein target
and mechanism do.} Unseen-drug generalisation is reachable through drug-side \emph{knowledge}, and
not through chemistry. Three caveats bound it. Coverage is $54\%$ and $32\%$, so this is measured
on the annotated subset. The channel gaps are severalfold below the advantage a same-drug lookup
enjoys for a \emph{seen} drug ($+0.44$ above chance), so the reachable ceiling in this regime is
much lower. And
target and MoA overlap, so whether target adds anything \emph{over} MoA is a further question this
measurement does not settle --- settling it would need a leave-one-mechanism-out split rather than
the leave-one-drug-out split used here, since a held-out drug's mechanism class is generally still
represented in training.

\subsection{The channel the model already has}
\label{sec:res-field}

The natural next step from the preceding section writes itself: append the protein targets to the
prompt, retrain, and see whether unseen-drug performance rises. That arm assumes the bottleneck is
\emph{information availability}. Everything in this thesis points the other way, but the direct
evidence is weaker than earlier drafts claimed, and this section exists to say exactly how weak it
is.

The scramble has until now replaced the drug name \emph{and} the mechanism together, so a non-null
gap establishes sensitivity to the combined prompt without attributing it to either span. On the
earlier target build, a decomposed version suggested the response is carried almost entirely by the
name span: swapping only the drug name gave $\gap = \ci{+0.0809}{+0.0592}{+0.1010}$ ($n = 570$),
swapping only the mechanism gave $+0.0091$ [$-0.0174$, $+0.0383$] ($n = 325$), and both together
$+0.1024$ [$+0.0827$, $+0.1236$]. Those arms were scored against the \texttt{opposite}-signature
comparator that Section~\ref{sec:res-generalisation} shows inflates every gap read against it, and
they were not re-run under the corrected evaluation. The claim that the model ``reads the name and
not the mechanism'' is therefore withdrawn pending that re-run; what the current run establishes is
response to the combined drug-and-mechanism prompt, with the mechanism field's contribution
unresolved.

The practical consequence is narrower than the one this section originally drew, and more honest.
The target-prompt arm is not retired by measurement --- it is untested. What argues for building it
carefully rather than eagerly is the shape of the failure documented above: the model's deficit is
not confined to missing information about unseen drugs, because on \emph{seen} drugs it recovers
only a fifth of the chance-to-ceiling distance against a lookup that knows nothing about contexts
at all. And what argues for building it at all is the one thing a lookup structurally cannot do:
handle a drug it has never seen, in the one regime where Section~\ref{sec:res-scope} shows live
channels exist. The model's opportunity is not the interaction --- Section~\ref{sec:res-transfer}
found $44\%$ of the residual variance there but no learnable structure in it --- it is the unseen
drug. Whether the model can be made to take that opportunity is the question the thesis leaves
open.
